\section{Gemma fixed guiding with full diagnostics}
\label{sec:gemma-cache-diagnostics}

The completed Gemma-4-26B-A4B NVFP4 experiment applies fixed positive
identity bias $+1$ on 1,395 held-out attempts from 465 problems, with
concurrency 16 and the frozen model-specific policy. Eight experts are
targeted at each of layers 16, 19, 21, 22, 25, 27 and 29. Calibration
uses replay routing from the corresponding Gemma checkpoint under a
different quantization; expert-identity transfer is therefore an assumption.
No held-out result is used to refit targets or choose the strength.

The original scorer gives 1,117 baseline and 1,113 guided correct answers.
Under the same offline extracted-answer equivalence contract used for
the other intervention evaluations, the counts are 1,223 and 1,222:
$87.67\%$ versus $87.60\%$. The change is $-0.07$ percentage points,
with a paired, dataset-stratified problem-bootstrap 95\% interval
$[-1.15,+1.00]$ (5,000 resamples, seed 42). There are 33 wrong-to-right
and 34 right-to-wrong changes. Mean generated tokens decrease from
11,695.48 to 11,562.61 ($-1.14\%$), and token-limit hits decrease from
67 to 66. These results show a small observed accuracy difference,
not established non-inferiority or an accuracy improvement.

The guided run was interrupted after 243 completed attempts and resumed
to full coverage. Answer-level analysis joins each completed ID exactly
once. Routing histograms sum the two saved execution sessions; they
include prefill, padding and potentially unfinished work from before the
interruption that was later regenerated. Consequently the diagnostic
population is execution work, not exactly one routing trace per scored
answer. Within each session the pre-bias and post-bias counters refer
to the same guided hidden states, not separate baseline generations.

\begin{table}[ht]
\centering\small
\caption{Gemma full-diagnostics fixed guiding: expert-selection coverage
in percent, aggregated across both execution sessions. Each static cache
contains eight experts per measured layer. Frequency pins use calibration
frequencies.}
\label{tab:gemma-cache-coverage}
\begin{tabular}{lrrr}
\toprule
Layer & Target, pre-bias & Target, post-bias & Frequency, pre-bias \\
\midrule
16 & 0.05 & 3.26 & 35.42 \\
19 & 0.03 & 1.06 & 39.41 \\
21 & 0.04 & 1.26 & 43.19 \\
22 & 0.11 & 2.02 & 53.53 \\
25 & 0.16 & 5.31 & 43.37 \\
27 & 1.38 & 6.42 & 41.37 \\
29 & 2.67 & 15.90 & 36.38 \\
All seven & 0.64 & 5.03 & 41.81 \\
\bottomrule
\end{tabular}
\end{table}

The target set's share of expert selections increases from $0.64\%$ to
$5.03\%$. For an already resident target cache this gives $4.43\%$ fewer
nonresident selections relative to the same targets before bias. However,
a same-budget frequency-selected cache covers $41.81\%$ before bias
(and $40.64\%$ after it). Guided target pinning consequently has
$63.20\%$ more nonresident selections than the pre-bias frequency cache.
The bias increases use of the selected identities but does not make them
a competitive static cache set for this model and strength.

These aggregate counters describe only seven policy layers and do not
measure ordered cache loads, cold-start cost, SSD traffic or latency.
They establish neither a whole-model speed effect nor a confidence interval
for cache coverage. The ordered experiments below test adaptive-cache
controls and all-layer loads separately.
